\section{Perancangan Solusi}
\subsection{Representasi Paket, Truk, dan State}
Sebagai notasi laporan, ruang peti kemas dinyatakan sebagai
$[0,W)\times[0,L)\times[0,H)$ dengan $z$ sebagai arah vertikal. Dimensi dasar
paket $i$ adalah $(w_i^0,\ell_i^0,h_i^0)$, beratnya $m_i$, nilainya $v_i$,
dan status fragile-nya $r_i\in\{0,1\}$. Orientasi $o_i$ memetakan dimensi dasar
menjadi dimensi efektif $(w_i,\ell_i,h_i)$.

Posisi sudut kiri-bawah-belakang paket dinotasikan sebagai
$\boldsymbol{p}_i=(x_i,y_i,z_i)$. Indikator $a_i(s)=1$ menyatakan paket $i$
dimuat dalam state $s$, sedangkan $a_i(s)=0$ menyatakan paket berada di luar.
Notasi indikator ini tidak mengharuskan implementasi memakai field boolean;
struktur data program tetap dapat dipilih bebas.

Representasi harus menyimpan identitas semua paket dan informasi ruang terisi
maupun kosong. Paket di luar tetap harus dapat ditampilkan, walaupun hanya
paket di dalam yang menyumbang nilai muatan. Dalam notasi laporan, keadaan
seluruh koordinat dapat diturunkan dari himpunan volume paket; implementasi
dapat memakai grid eksplisit atau representasi lain yang ekuivalen.


\subsection{Pembentukan Initial State}
Spesifikasi mengharuskan inisialisasi acak. Prosedur pembentukan initial state
perlu menjelaskan sumber keacakan pada pemilihan paket, posisi, atau orientasi,
serta cara menjamin atau memulihkan validitas konfigurasi. Seed dicatat agar
run dapat direproduksi. Pada GA, pencatatan diperluas ke populasi awal.


\subsection{Pembentukan Neighbor}
\subsubsection{Pertukaran Dua Paket}
Operator pertukaran melibatkan dua paket, termasuk kombinasi paket di dalam
dan di luar truk. Karena dimensi dan orientasi dapat berbeda, hasil pertukaran
harus diperiksa kembali terhadap batas ruang, tabrakan, support, fragile,
dan kapasitas berat. Pemindahan paket penopang juga dapat mengubah validitas
paket lain yang tidak dipilih langsung.

\subsubsection{Pemindahan Sebuah Paket}
Operator pemindahan mengubah koordinat suatu paket. Berdasarkan QnA, operator
ini dapat dipakai untuk memasukkan paket ke truk atau mengeluarkannya.
Representasi tujuan di luar truk perlu mempunyai aturan yang jelas dan tidak
menghasilkan paket yang hanya sebagian berada di dalam.

\subsubsection{Rotasi Sebuah Paket}
Operator rotasi mengubah orientasi sebesar $90^\circ$ terhadap salah satu
sumbu. Pertukaran dimensi efektif masing-masing dapat ditulis sebagai
$(w,\ell,h)\mapsto(w,h,\ell)$ untuk sumbu $x$,
$(w,\ell,h)\mapsto(h,\ell,w)$ untuk sumbu $y$, dan
$(w,\ell,h)\mapsto(\ell,w,h)$ untuk sumbu $z$. Notasi ini menjelaskan ukuran
efektif; lokasi titik acuan setelah rotasi tetap perlu didefinisikan program.


\subsection{Fungsi Objektif}
Untuk himpunan state valid $\mathcal{F}$, persoalan wajib diformulasikan sebagai
\begin{equation}
 \max_{s\in\mathcal{F}} f(s),\qquad
 f(s)=\sum_{i=1}^{N}a_i(s)v_i.\label{eq:objektif}
\end{equation}
Nilai paket di luar tidak dihitung sebagai nilai muatan. Jika program memakai
penalti bagi state tidak valid atau heuristik untuk memilih kandidat, fungsi
tersebut harus dibedakan dari $f(s)$. Nilai objektif yang dilaporkan untuk
solusi akhir tetap mengikuti Persamaan~\eqref{eq:objektif} pada state valid.

\subsection{Pemeriksaan Validitas State}\label{sec:validasi}
\subsubsection{Batas Ruang dan Tumpang Tindih}
Untuk setiap paket yang dimuat, berlaku
\begin{equation}
 0\leq x_i,\quad x_i+w_i\leq W,\quad
 0\leq y_i,\quad y_i+\ell_i\leq L,\quad
 0\leq z_i,\quad z_i+h_i\leq H.
\end{equation}
Seluruh koordinat posisi berupa bilangan bulat. Dengan volume setengah
terbuka $V_i=[x_i,x_i+w_i)\times[y_i,y_i+\ell_i)\times[z_i,z_i+h_i)$,
dua paket yang dimuat tidak boleh memiliki irisan volume:
$V_i\cap V_j=\varnothing$ untuk $i\ne j$. Kontak pada permukaan tidak
dianggap tumpang tindih volume.

\subsubsection{Support dan Paket Melayang}
Misalkan proyeksi alas paket $i$ adalah
$A_i=[x_i,x_i+w_i)\times[y_i,y_i+\ell_i)$.
Paket $j$ menjadi penopang geometris paket $i$ apabila
\begin{equation}
 j\in S_i(s)\iff a_j(s)=1,\ j\ne i,\ z_j+h_j=z_i,
 \quad A_i\cap A_j\ne\varnothing.
\end{equation}
Untuk setiap paket di dalam, syarat tidak melayang adalah
\begin{equation}
 z_i=0\quad\lor\quad S_i(s)\ne\varnothing.\label{eq:support}
\end{equation}
Syarat ini mengikuti definisi eksistensi support. Dukungan sebagian alas
diperbolehkan. Sentuhan tepi tanpa irisan pada domain alas setengah terbuka
tidak dihitung sebagai area penopang. Prosedur pemeriksaan diskret harus
konsisten dengan konvensi koordinat yang digunakan.

\subsubsection{Larangan Fragile Menjadi Penopang}
Larangan fragile dinyatakan secara terpisah dari kecukupan support:
\begin{equation}
 \forall i\text{ yang dimuat},\quad \forall j\in S_i(s),\quad r_j=0.
\end{equation}
Artinya, keberadaan penopang non-fragile tidak membenarkan kontak support lain
yang berasal dari paket fragile. Paket fragile sendiri tetap boleh diletakkan
di lantai atau di atas paket non-fragile, selama tidak menopang paket lain.

\subsubsection{Kapasitas Berat}
Jika kapasitas maksimum truk adalah $C$, kendala muatan adalah
\begin{equation}
 \sum_{i=1}^{N}a_i(s)m_i\leq C.\label{eq:berat}
\end{equation}
Kapasitas ini tidak menyatakan jumlah paket atau volume. Kendala ruang dan
berat diperiksa secara terpisah.


\subsection{Rancangan Algoritma Local Search}
Local search mengevaluasi konfigurasi kandidat untuk mencari state bernilai
tinggi. Tiga pendekatan pada tugas ini berbeda dalam pemilihan kandidat dan
mekanisme keluar dari konfigurasi yang sulit diperbaiki~\cite{kuliah}.

\subsubsection{Hill-Climbing}
Hill-climbing berusaha meningkatkan nilai state melalui neighbor. Pada
steepest-ascent, neighbor terbaik dipilih dari neighborhood yang dievaluasi.
Pada stochastic hill-climbing dalam salindia, kandidat dibangkitkan acak dan
perpindahan dilakukan jika nilainya lebih baik. Sideways move mengizinkan
langkah dengan nilai sama sesuai batas yang ditetapkan, sedangkan random
restart mengulangi pencarian dari initial state acak~\cite{kuliah}.


\subsubsection{Simulated Annealing}
SA dapat menerima langkah yang menurunkan nilai objektif. Dengan
$\Delta E=f(s')-f(s)$ dan temperatur $T>0$, probabilitas penerimaan dirumuskan
sebagai~\cite{kuliah}
\begin{equation}
 p_{\mathrm{terima}}=
 \begin{cases}1,&\Delta E\geq0,\\
 \exp(\Delta E/T),&\Delta E<0.
 \end{cases}\label{eq:sa}
\end{equation}
Nilai $\exp(\Delta E/T)$ untuk kandidat membaik dapat melebihi satu; nilai
tersebut berbeda dari probabilitas penerimaan yang dibatasi satu. Grafik
eksperimen harus menjelaskan apakah menampilkan ekspresi yang diminta
spesifikasi atau probabilitas penerimaan, serta domain kandidatnya. Untuk
$T=0$, pencarian dihentikan sebelum menghitung pembagian.


\subsubsection{Genetic Algorithm}
GA memelihara populasi individu yang merepresentasikan kandidat solusi.
Seleksi memilih induk, crossover menggabungkan informasi induk, dan mutasi
memperkenalkan perubahan. Representasi harus memungkinkan hasil operator
diterjemahkan kembali menjadi konfigurasi paket yang dapat divalidasi~\cite{kuliah}.


\subsection{Format Input dan Syarat Validitas}
Dokumentasi input harus menyebut ekstensi/struktur berkas, nama field, tipe
data, satuan, orientasi, serta representasi posisi apabila posisi diberikan
sebagai input. Validasi perlu mencakup identitas unik, dimensi positif,
kapasitas, domain atribut numerik, dan nilai boolean fragile.
